\begin{textsample}{POS Dim 7 – human – Score 4.00 – es/es\_ef\_anos\_iniciais\_matematica\_3\_p15.txt}  \label{ex:f7_pos_005}
Também é importante ressaltar que as habilidades \textbf{representam} o que se \textbf{espera} que os estudantes aprendam ao longo de cada ano do Ensino Fundamental e não \textbf{descrevem} ações ou condutas docentes, nem induzem à opção por \textbf{abordagens} ou metodologias, que devem ser adotadas, adequando -se à realidade de cada unidade de ensino, considerando o contexto e as características de seus estudantes ( BRASIL, 2017 ).
No Currículo de Matemática do Espírito Santo, com as contribuições, foi -se necessário alterar as terminologias de algumas habilidades, passando a representação para EF01MA01/ES.
Destaca -se também a inserção de novas habilidades, sem alteração da ordenação das habilidades mínimas da BNCC e com a garantia de progressão e especificidade \textbf{regional}.

% matched lemmas: abordagem, descrever, esperar, regional, representar
\end{textsample}
